
%\documentclass[10pt,twoside,twocolumn]{article}
\documentclass[12pt,twoside]{article}
\usepackage[bf,small]{caption}
\usepackage[letterpaper,hmargin=1in,vmargin=1in]{geometry}
\usepackage{paralist} % comapctitem, compactdesc, compactenum
\usepackage{titlesec}
\usepackage{titletoc}
\usepackage{times}
\usepackage{hyperref}
\usepackage{algorithmic}
\usepackage{graphicx}
\graphicspath{{./graphics/}}
\usepackage{xspace}
\usepackage{verbatim}
\usepackage{url}
\usepackage{float}
\hyphenation{Sub-Bytes Shift-Rows Mix-Col-umns Add-Round-Key}

\setlength{\parskip}{12pt}
\setlength{\parindent}{0pt}

\newcommand{\hdb}{\emph{hashdb}\xspace}
\newcommand{\bulk}{\emph{bulk\_extractor}\xspace}
\newcommand{\hashid}{\emph{hashid}\xspace}
\newcommand{\mdd}{\emph{md5deep}\xspace}
\newcommand{\bev}{\emph{Bulk Extractor Viewer}\xspace}

\begin{document}

\begin{center}
\Large Demo: Finding Part of a Video File in a Media Image \\
\large Using \bulk and \hdb
\end{center}

In this demo, we find that a media image contains part of a
``blacklisted'' video file.
This demo uses the following resources:
\begin{compactitem}
\item The media image containing a fragment of a video file.
\item A \hdb block hash database containing block hashes
from the blacklist video file.
\item \hdb.
\item \bulk compiled with the \hdb \hashid scanner.
\end{compactitem}

Here is the workflow:

\begin{figure}[H]
  \center
  \includegraphics[scale=0.6]{drawings/scan_hashid}
  \caption*{Scan the media image for parts of a video file.}
  \label{fig:scan_hashid}
\end{figure}

Steps:
\begin{compactenum}
\item Download and install \hdb from
\url{http://digitalcorpora.org/downloads/hashdb}
as described at
\url{https://github.com/simsong/hashdb/wiki/Installing-hashdb}.
  \begin{compactitem}
  \item \textbf{Windows}: Download and install the latest hashdb windowsinstaller.exe file:
  \url{http://digitalcorpora.org/downloads/hashdb/hashdb-1.0.0-beta-1-windowsinstaller.exe}
  \item \textbf{Linux, Mac}: Download and build the latest hashdb distribution:
  \url{http://digitalcorpora.org/downloads/hashdb/hashdb-1.0.0-beta-1.tar.gz}.
  \end{compactitem}
\item Download and install \bulk compiled with \hdb from
\url{http://digitalcorpora.org/downloads/hashdb}
as described at
\url{https://github.com/simsong/hashdb/wiki/Installing-hashdb}.

Currently, \bulk compiled with \hdb
is not available from the \bulk distribution.
I am providing the custom builds shown here until the next release of \bulk.
  \begin{compactitem}
  \item \textbf{Windows}: Download and install the latest
  bulk\_extractor windowsinstaller.exe file
  compiled with the \hdb \hashid scanner:
  \url{http://digitalcorpora.org/downloads/hashdb/bulk\_extractor-1.4.1-windowsinstaller.exe}
  \item \textbf{Linux, Mac}: Download and build the latest \bulk:
  \url{http://digitalcorpora.org/downloads/hashdb/bulk\_extractor-1.4.5devel.tar.gz}.
  \end{compactitem}

\item Download the data files required for the demo into a working directory:
  \begin{compactitem}
  \item The database containing hash blocks of the demo video file:
  \url{http://digitalcorpora.org/downloads/hashdb/demo/demo\_video\_db.zip}
  \item Media containing a fragment of the demo video file,
  specifically, a contiguous 64KiB section
  near the end of about 10 MiB of video data:
  \url{http://digitalcorpora.org/downloads/hashdb/demo/demo\_video\_redacted\_image}
  \end{compactitem}

\item Extract the demo database
by unzipig the \texttt{demo\_video\_db.zip} file
into your working directory.
Please check that it doesn't unzip into a nested subfolder.
\item Scan for matching hash values: Using a command window, go to your working directory and then run \bulk,
specifying the paths to the hash database and the media:
\begin{verbatim}
$ bulk_extractor -o outdir -S hashdb_mode=scan \
  -S hashdb_path_or_socket=demo_video_db \
  demo_video_redacted_image
\end{verbatim}
\item View the feature file using an editor or use the \bev tool.
For example to view with Windows Notepad, type:
\begin{verbatim}
$ notepad outdir/identified_blocks.txt
\end{verbatim}
An example hash block match looks like this:
\begin{verbatim}
12452352    3b6b477d391f73f67c1c01e2141dbb17    1
\end{verbatim}
\end{compactenum}

Seeing hash \texttt{3b6b477...} at Forensic path \texttt{12452352}
shows that a hash block match was found, but what file does it match?
We find the file that contains the hash by using a \hdb source lookup:
\begin{figure}[H]
  \center
  \includegraphics[scale=0.6]{drawings/source_lookup}
  \caption*{Look up the file that has the hash.}
  \label{fig:source_lookup}
\end{figure}

Steps to look up source information about the identified blocks:
\begin{compactenum}
\item Using a command window, go to your working directory and then run
the \hdb tool:
\begin{verbatim}
$ hashdb expand_identified_blocks demo_video_db \
  outdir/identified_blocks.txt > outdir/identified_sources.txt
\end{verbatim}

\item Now view file \texttt{outdir/identified\_sources.txt} to see
features containing source information.
This example line:
\begin{verbatim}
12452352    3b6b477d391f73f67c1c01e2141dbb17 \
repository_name=repository_demo_video_dfxml, \
filename=/home/bdallen/demo/demo_video.mp4, \
file_offset=10485760
\end{verbatim}

states that the block at Forensic path \texttt{12452352}
matches the block \texttt{10485760} bytes into the
\texttt{demo\_video.mp4} video file
in the hash database,
indicating a positive match.
\end{compactenum}

This completes the demo.

\end{document}

